\chapter{Примеры тестовых данных и результатов}
\label{app:examples}

В данном приложении приведены примеры входных данных и результатов работы алгоритма.

\section{Пример 1: Однородная укладка}

\subsection{Входные данные}
6 коробок одного типа размером $600 \times 400 \times 200$ мм.

\begin{verbatim}
SKU,Quantity,Length,Width,Height,Weight,Strength,Aisle,Caustic
900001,6,600,400,200,12000,4,1,0
\end{verbatim}

\subsection{Результат}
\begin{verbatim}
SKU,x,y,z,X,Y,Z
900001,0,0,0,600,400,200
900001,600,0,0,1200,400,200
900001,0,400,0,600,800,200
900001,600,400,0,1200,800,200
900001,0,0,200,600,400,400
900001,600,0,200,1200,400,400
\end{verbatim}

\textbf{Метрики:}
\begin{itemize}
    \item Высота паллеты: 400 мм
    \item Percolation: 75.0\%
\end{itemize}

% TODO: Добавить визуализацию примера 1 (figures/example1_uniform.png)
% Показать: 6 одинаковых коробок в 2 слоя по 4 штуки (2x2)

\section{Пример 2: Смешанная укладка}

\subsection{Входные данные}
Коробки разных размеров:

\begin{verbatim}
SKU,Quantity,Length,Width,Height,Weight,Strength,Aisle,Caustic
900001,4,600,400,200,12000,4,1,0
900002,6,200,200,100,2000,4,2,0
900003,2,400,300,150,5000,4,1,0
\end{verbatim}

\subsection{Характеристики}
\begin{itemize}
    \item Всего коробок: 12
    \item Суммарный объём: $4 \times 48000000 + 6 \times 4000000 + 2 \times 18000000 = 252000000$ мм³
    \item Размер паллеты: $1200 \times 800$ мм
\end{itemize}

% TODO: Добавить визуализацию примера 2 (figures/example2_mixed.png)
% Показать: разноцветные коробки разных размеров

\section{Пример 3: Реальный тест (tests/100.csv)}

\subsection{Характеристики входных данных}
\begin{itemize}
    \item Количество типов коробок: 15
    \item Общее количество коробок: 47
    \item Минимальный размер: $100 \times 80 \times 50$ мм
    \item Максимальный размер: $500 \times 400 \times 300$ мм
\end{itemize}

\subsection{Результаты алгоритмов}
\begin{table}[ht]
\centering
\begin{tabular}{|l|c|c|c|}
\hline
\textbf{Алгоритм} & \textbf{Высота} & \textbf{Percolation} & \textbf{Время} \\
\hline
GreedySolver & 1250 мм & 71.2\% & 2 мс \\
LNSSolver (1 сек) & 1180 мм & 75.4\% & 1.0 сек \\
LNSSolver (30 сек) & 1120 мм & 79.6\% & 30 сек \\
\hline
\end{tabular}
\end{table}

% TODO: Добавить 3 визуализации рядом (figures/example3_comparison.png)
% Показать: результаты GreedySolver, LNS(1с), LNS(30с) для сравнения

\section{Анализ распределения результатов}

На основе 436 тестов построено распределение коэффициента заполнения:

% TODO: Добавить гистограмму распределения percolation (figures/percolation_histogram.png)
% Показать: гистограмма с двумя сериями (GreedySolver и LNSSolver)
% Ось X: percolation (50%-100%), Ось Y: количество тестов

\subsection{Статистика по тестовым данным}

\begin{table}[ht]
\centering
\caption{Статистика тестового набора}
\begin{tabular}{|l|c|}
\hline
\textbf{Параметр} & \textbf{Значение} \\
\hline
Количество тестов & 436 \\
Среднее число коробок в тесте & 63 \\
Минимальное число коробок & 5 \\
Максимальное число коробок & 312 \\
Средний объём коробки & 12.4 литра \\
\hline
\end{tabular}
\end{table}

\section{Типичные паттерны укладки}

\subsection{Слоистая укладка}
Возникает когда коробки имеют одинаковую высоту. Характеризуется высоким stability (0.2--0.4), но низким interlocking.

% TODO: Добавить пример слоистой укладки (figures/pattern_layered.png)

\subsection{Кирпичная укладка}
Коробки смещены относительно нижнего слоя. Высокий interlocking (0.6--0.9), хорошая устойчивость.

% TODO: Добавить пример кирпичной укладки (figures/pattern_brick.png)

\subsection{Смешанная укладка}
Коробки разных размеров заполняют пустоты. Типичный случай для реальных данных.

% TODO: Добавить пример смешанной укладки (figures/pattern_mixed.png)
